% lecture.tex
\documentclass[aspectratio=169, xcolor={table,xcdraw}]{beamer}

\usetheme{Madrid}
\usecolortheme{default}
\definecolor{ucbblue}{RGB}{0, 51, 102}

\setbeamercolor{structure}{fg=ucbblue}
\setbeamercolor*{palette primary}{fg=white,bg=ucbblue}
\setbeamercolor*{palette secondary}{fg=white,bg=ucbblue!80}
\setbeamercolor*{palette tertiary}{fg=white,bg=ucbblue!90}
\setbeamercolor{title}{fg=white, bg=ucbblue}
\setbeamercolor{frametitle}{fg=white, bg=ucbblue}
\setbeamercolor{item}{fg=ucbblue}
\setbeamercolor{block title}{bg=ucbblue!15, fg=ucbblue}
\setbeamercolor{block body}{bg=ucbblue!5}

\usepackage[T1]{fontenc}
\usepackage[utf8]{inputenc}
\usepackage{tgpagella}
\usefonttheme{serif}
\usepackage[spanish, es-noshorthands]{babel}
\usepackage{amsmath, amssymb, bm}
\usepackage{graphicx}
\usepackage[most]{tcolorbox}
\usepackage{booktabs}
\usepackage{tabularx}
\usepackage{tikz}
\usepackage{pgfplots}
\pgfplotsset{compat=1.18}
\usepackage[american]{circuitikz}
\usetikzlibrary{shapes.geometric, arrows.meta, positioning, calc, backgrounds}

\title[Sensor y Planta Térmica]{Calibración de Sensor e Identificación de Planta Térmica}
\subtitle{IMT-221 Circuitos Electrónicos II — Sesión 2/10 (Tutoría Intensiva)}
\author[B. Quiroga Turdera]{M.Sc. Ing. Bernardo Quiroga Turdera}
\institute[UCB Tarija]{Universidad Católica Boliviana ``San Pablo'' — Sede Tarija}
\date{Sesión 2}

\begin{document}

\begin{frame}
    \titlepage
\end{frame}

\begin{frame}{Pregunta de Ingeniería}
    \begin{tcolorbox}[colback=ucbblue!5, colframe=ucbblue, title=Objetivo de la Sesión]
        \footnotesize Antes de diseñar un acondicionador (op-amp) o un controlador (PI/PID), necesitamos caracterizar dos fenómenos independientes: la \textbf{calibración del sensor} y la \textbf{dinámica de la planta}.
    \end{tcolorbox}
    
    \vspace{0.3cm}
    \centering
    \textbf{V3: Arquitectura de Caracterización (Dos Modelos Separados)}
    
    \vspace{0.4cm}
    \begin{tikzpicture}[node distance=0.6cm and 0.6cm, thick,
        block/.style={draw, rounded corners, fill=blue!10, text width=2.2cm, align=center, minimum height=0.9cm, font=\footnotesize}]
        
        % Modelo 1: Calibración
        \node[block, fill=green!10] (ref) {Ref. Térmica\\$T_{ref}$};
        \node[block, right=of ref] (calib) {Calibración\\$V_s = a T_{ref} + b$};
        \node[block, fill=orange!10, right=of calib] (lectura) {Lectura\\$V_{s}$};
        
        \draw[->] (ref) -- (calib);
        \draw[->] (lectura) -- (calib);
        
        % Modelo 2: Dinámica
        \node[block, fill=red!10, below=0.8cm of ref] (actuador) {Comando\\$u(t)$};
        \node[block, right=of actuador] (planta) {Planta\\$G_p(s)$};
        \node[block, fill=orange!10, right=of planta] (sensor) {Sensor\\$T(t) \to V_s(t)$};
        
        \draw[->] (actuador) -- (planta);
        \draw[->] (planta) -- (sensor);
    \end{tikzpicture}
\end{frame}

\begin{frame}{Recuperación CP1A: RC Pasa-Bajos}
    \begin{columns}[T]
        \begin{column}{0.48\textwidth}
            \centering
            \textbf{Topología (Filtro RC)}
            
            \vspace{0.2cm}
            \begin{circuitikz}[thick, scale=0.8, transform shape]
                \draw (0,2) to[short, o-] (1,2) to[R, l=$R$] (3,2) to[short, -o] (4.5,2);
                \draw (3,2) to[C, l=$C$] (3,0);
                \draw (0,0) to[short, o-o] (4.5,0);
                \draw (3,0) node[ground]{};
                
                \node[left] at (0,2) {$v_{in}$};
                \node[right] at (4.5,2) {$v_{out}$};
            \end{circuitikz}
            
            \vspace{0.4cm}
            \textbf{Normalización y Corte:}
            
            \vspace{0.2cm}
            {\centering $\displaystyle \omega_c = \frac{1}{RC} \implies f_c = \frac{1}{2\pi RC}$ \par}
            \vspace{0.2cm}
            {\centering $\displaystyle H(j\omega) = \frac{1}{1 + j\left(\frac{f}{f_c}\right)}$ \par}
        \end{column}
        \begin{column}{0.48\textwidth}
            \textbf{Funciones Analíticas:} 
            
            \vspace{0.2cm}
            {\centering $\displaystyle G_{dB}(f) = -10 \log_{10}\left[1 + \left(\frac{f}{f_c}\right)^2\right]$ \par}
            \vspace{0.2cm}
            {\centering $\displaystyle \phi(f) = -\arctan\left(\frac{f}{f_c}\right)$ \par}
            
            \vspace{0.4cm}
            \begin{tcolorbox}[colback=white, colframe=gray, title=Ejemplo Ilustrativo (No Medido)]
                \scriptsize Si $R=10\,\text{k}\Omega, C=100\,\text{nF} \implies f_c \approx 159.1\,\text{Hz}$
                \begin{itemize}
                    \item En $0.1 f_c$: $G_{dB} \approx -0.04\,\text{dB}$, $\phi \approx -5.7^\circ$
                    \item En $f_c$: $G_{dB} \approx -3.01\,\text{dB}$, $\phi = -45^\circ$
                    \item En $10 f_c$: $G_{dB} \approx -20.0\,\text{dB}$, $\phi \approx -84.3^\circ$
                \end{itemize}
            \end{tcolorbox}
        \end{column}
    \end{columns}
\end{frame}

\begin{frame}{V4: Inventario y Reglas de Seguridad}
    \begin{alertblock}{Restricción de Energización}
        \footnotesize No existe permiso técnico para energizar una planta térmica solo porque esta guía lo describe. Si las conexiones y límites no están verificados, se utilizará un \textbf{Dataset Histórico Auténtico} provisto por el instructor.
    \end{alertblock}
    
    \vspace{0.2cm}
    \begin{center}
    \resizebox{0.95\textwidth}{!}{
    \begin{tabular}{@{} l l l @{}}
        \toprule
        \textbf{Requisito de Inventario} & \textbf{Estado} & \textbf{Procedencia de los Datos} \\ \midrule
        Actuador y entrada manipulada $u(t)$ & CONFIRMADO / FALTA & (Manual, Ficha, Inspección) \\ \addlinespace
        Sensor y señal cruda observada $V_s$ & CONFIRMADO / FALTA & (Multímetro, Ficha Técnica) \\ \addlinespace
        Instrumento patrón de referencia $T_{ref}$ & CONFIRMADO / FALTA & (Termocupla validada, Termómetro) \\ \addlinespace
        Alimentación y límites seguros & CONFIRMADO / FALTA & (Fuentes de laboratorio) \\ \addlinespace
        Protocolo de apagado de emergencia & CONFIRMADO / FALTA & (Procedimiento UCB) \\ \bottomrule
    \end{tabular}
    }
    \end{center}
\end{frame}

\begin{frame}{Calibración Estática del Sensor}
    \begin{columns}[T]
        \begin{column}{0.48\textwidth}
            La calibración busca una relación empírica entre la lectura $V_s$ y la referencia $T_{ref}$:
            
            \vspace{0.2cm}
            {\centering $\displaystyle V_s = a T_{ref} + b + \epsilon$ \par}
            
            \vspace{0.4cm}
            Minimizando el error cuadrático:
            
            \vspace{0.2cm}
            {\centering $\displaystyle a = \frac{\sum (T_i - \bar{T})(V_i - \bar{V})}{\sum (T_i - \bar{T})^2} \qquad b = \bar{V} - a\bar{T}$ \par}
            
            \vspace{0.4cm}
            \footnotesize \textbf{Nota:} Invertir una regresión de $V$ sobre $T$ no es idéntico a ajustar $T$ sobre $V$.
        \end{column}
        \begin{column}{0.48\textwidth}
            \centering
            \textbf{V5: Ejemplo Ilustrativo de Ajuste}
            
            \begin{tikzpicture}
                \begin{axis}[
                    width=0.95\linewidth, height=4.8cm, % Se elimina scale only axis
                    xlabel={$T_{ref}$ ($^\circ$C)}, ylabel={$V_s$ (V)},
                    grid=both,
                    xmin=15, xmax=55, ymin=0.4, ymax=1.2,
                    label style={font=\small}, tick label style={font=\footnotesize},
                    legend style={at={(0.05,0.95)}, anchor=north west, font=\tiny}
                ]
                % Recta ajustada V = 0.02*T + 0.1
                \addplot[thick, blue, domain=20:50] {0.02*x + 0.1};
                \addlegendentry{Ajuste: $0.02T + 0.1$}
                % Puntos "medidos"
                \addplot[only marks, mark=*, red] coordinates {
                    (20, 0.51) (30, 0.68) (40, 0.91) (50, 1.10)
                };
                \addlegendentry{Pares $(T_i, V_i)$}
                \end{axis}
            \end{tikzpicture}
        \end{column}
    \end{columns}
\end{frame}

\begin{frame}{Identificación Dinámica: Ensayo al Escalón}
    Aplicamos un cambio conocido $\Delta u$ al actuador y observamos las variables en desviación: $\delta u(t) = u(t) - u_0$, $\delta T(t) = T(t) - T_0$.
    
    \vspace{0.3cm}
    Para un modelo continuo de \textbf{primer orden sin retardo}:
    
    \vspace{0.2cm}
    {\centering $\displaystyle \tau \frac{d\delta T(t)}{dt} + \delta T(t) = K \delta u(t) \implies G_p(s) = \frac{\Delta T(s)}{\Delta U(s)} = \frac{K}{\tau s + 1}$ \par}
    
    \vspace{0.3cm}
    Si observamos una estabilización completa ($T_\infty$):
    
    \vspace{0.2cm}
    {\centering $\displaystyle K = \frac{T_\infty - T_0}{\Delta u} \quad \left[ \frac{^\circ\text{C}}{\text{Unidad}} \right]$ \par}
    
    \vspace{0.3cm}
    El parámetro temporal $\tau$ se extrae en el instante donde la respuesta alcanza el $63.2\%$ de su recorrido total:
    
    \vspace{0.2cm}
    {\centering $\displaystyle T_{63} = T_0 + 0.632(T_\infty - T_0)$ \par}
\end{frame}

\begin{frame}{V6: Análisis Gráfico del Escalón (Ilustrativo)}
    \begin{columns}[T]
        \begin{column}{0.48\textwidth}
            \centering
            \textbf{Comando $u(t)$}
            
            \begin{tikzpicture}
                \begin{axis}[
                    width=0.95\linewidth, height=4.2cm, % Se elimina scale only axis
                    xlabel={$t$ (s)}, ylabel={$u(t)$},
                    grid=both,
                    xmin=0, xmax=100, ymin=0, ymax=0.6,
                    label style={font=\small}, tick label style={font=\footnotesize}
                ]
                \addplot[thick, blue, const plot] coordinates {(0,0.2) (10,0.2) (10,0.5) (100,0.5)};
                \node[above left, font=\scriptsize] at (axis cs: 10,0.2) {$u_0=0.2$};
                \node[above left, font=\scriptsize] at (axis cs: 50,0.5) {$\Delta u=0.3$};
                \end{axis}
            \end{tikzpicture}
        \end{column}
        \begin{column}{0.48\textwidth}
            \centering
            \textbf{Respuesta Térmica $T(t)$}
            
            \begin{tikzpicture}
                \begin{axis}[
                    width=0.95\linewidth, height=4.2cm, % Se elimina scale only axis
                    xlabel={$t$ (s)}, ylabel={$T(t)$ ($^\circ$C)},
                    grid=both,
                    xmin=0, xmax=100, ymin=20, ymax=35,
                    label style={font=\small}, tick label style={font=\footnotesize}
                ]
                \addplot[thick, red, domain=0:10] {25};
                \addplot[thick, red, domain=10:100, samples=50] {25 + 6*(1 - exp(-(x-10)/40))};
                % Marcas
                \draw[dashed, black] (axis cs: 0, 31) -- (axis cs: 100, 31) node[pos=0.1, above, font=\scriptsize]{$T_\infty=31$};
                \draw[dashed, black] (axis cs: 50, 25) -- (axis cs: 50, 28.79) -- (axis cs: 0, 28.79) node[left, font=\scriptsize]{$T_{63}$};
                \end{axis}
            \end{tikzpicture}
        \end{column}
    \end{columns}
    
    \vspace{0.2cm}
    {\centering \footnotesize $\displaystyle K = \frac{31 - 25}{0.3} = 20\, \frac{^\circ\text{C}}{\text{Unidad}} \qquad \tau = 50 - 10 = 40\,\text{s} \qquad \implies G_p(s) = \frac{20}{40s + 1}$ \par}
\end{frame}

\begin{frame}{Validación y Trazabilidad del Modelo}
    Un ajuste empírico es solo una hipótesis. Se debe contrastar frente a observaciones reales mediante el análisis de residuos $e_i$:
    
    \vspace{0.2cm}
    {\centering $\displaystyle e_i = T_{obs,i} - T_{mod}(t_i)$ \par}
    \vspace{0.2cm}
    {\centering $\displaystyle RMSE_T = \sqrt{\frac{1}{N}\sum e_i^2} \quad \quad \quad MAE_T = \frac{1}{N}\sum \vert{}e_i\vert{}$ \par}
    
    \vspace{0.4cm}
    \begin{alertblock}{Diagnóstico Experimental}
        \footnotesize
        \begin{itemize}
            \item \textbf{Inicio Tardío:} Sugiere un retardo físico $L \implies G_{p,L}(s) = \frac{K e^{-Ls}}{\tau s + 1}$.
            \item \textbf{Error Creciente:} Posible deriva ambiental o falla en alcanzar estacionario.
            \item \textbf{Trazabilidad:} Datos sin sello de tiempo o sin fuente declarada \textbf{carecen de validez técnica} y quedan en estado "PENDIENTE".
        \end{itemize}
    \end{alertblock}
\end{frame}

\begin{frame}{V5: Diagrama de Bode Analítico Ideal}
    \centering
    \begin{tikzpicture}
        \begin{semilogxaxis}[
            name=mag,
            width=0.85\textwidth, height=3.2cm, % Reducido y sin scale only axis
            ylabel={$G_{dB}$ (dB)},
            xticklabel=\empty,
            ymin=-25, ymax=5,
            ytick={-20, -3, 0},
            grid=both,
            domain=0.01:100,
            samples=100,
            label style={font=\small},
            tick label style={font=\footnotesize}
        ]
            \addplot[thick, blue] {-10*log10(1 + x^2)};
            \addplot[dashed, red] coordinates {(0.01, 0) (1, 0)};
            \addplot[dashed, red] { -20*log10(x) };
            \addplot[mark=*, black, mark size=1.5pt] coordinates {(1, -3.01)} node[below left, font=\scriptsize] {$-3\text{ dB}$};
        \end{semilogxaxis}
        
        \begin{semilogxaxis}[
            at={(mag.south west)},
            anchor=north west,
            yshift=-0.1cm, % Acercado para optimizar espacio vertical
            width=0.85\textwidth, height=3.2cm,
            xlabel={$f/f_c$ (Normalizada)},
            ylabel={$\phi$ ($^\circ$)},
            ymin=-95, ymax=5,
            ytick={0, -45, -90},
            grid=both,
            domain=0.01:100,
            samples=100,
            label style={font=\small},
            tick label style={font=\footnotesize}
        ]
            \addplot[thick, blue] {-atan(x)};
            \addplot[mark=*, black, mark size=1.5pt] coordinates {(1, -45)} node[above right, font=\scriptsize] {$-45^\circ$};
        \end{semilogxaxis}
    \end{tikzpicture}
\end{frame}

\begin{frame}{Requisitos de Acondicionamiento (Para Sesión 3)}
    \begin{tcolorbox}[colback=white, colframe=ucbblue, title=V8: Especificaciones de Rango]
        Transformamos el rango de salida del sensor $[V_{s,min}, V_{s,max}]$ a un rango objetivo del controlador $[V_{o,min}, V_{o,max}]$ mediante una relación afín.
    \end{tcolorbox}
    
    \vspace{0.2cm}
    La ecuación de diseño general para el op-amp será: 
    
    \vspace{0.2cm}
    {\centering $\displaystyle V_o = A V_s + B$ \par}
    
    \vspace{0.4cm}
    {\centering $\displaystyle A = \frac{V_{o,max} - V_{o,min}}{V_{s,max} - V_{s,min}} \qquad \text{y} \qquad B = V_{o,min} - A V_{s,min}$ \par}
    
    \vspace{0.4cm}
    \begin{alertblock}{Nota Crítica}
        \scriptsize Un cálculo matemático de $A$ y $B$ no garantiza que un amplificador real pueda generarlos. Falta revisar \textit{Output Swing}, \textit{Common-Mode} y alimentación (Sesión 3).
    \end{alertblock}
\end{frame}

\begin{frame}[fragile]{Metodología de Validación: Simulación y Laboratorio}
    \begin{columns}[T]
        \begin{column}{0.4\textwidth}
            \textbf{Starter SPICE (.ac):}
\begin{verbatim}
* CP1A - RC low pass 
V1 in 0 AC 1
R1 in out 10k
C1 out 0 100n
.ac dec 100 10 10000
.end
\end{verbatim}
            \vspace{0.2cm}
            \small \textit{Graficar $V(out)/V(in)$ para observar la transferencia $H(s)$ independientemente de la fuente física.}
        \end{column}
        \begin{column}{0.58\textwidth}
            \centering
            \textbf{V6: Flujo de Contraste y Discrepancia}\\
            \vspace{0.2cm}
            \begin{tikzpicture}[node distance=0.8cm, thick,
                block/.style={draw, rounded corners, fill=blue!10, text width=2cm, align=center, minimum height=1.1cm, font=\scriptsize}]
                
                \node[block] (pred) {\textbf{Predicción}\\Ecuación};
                \node[block, right=0.3cm of pred] (sim) {\textbf{Simulación}\\Barrido AC};
                \node[block, right=0.3cm of sim] (med) {\textbf{Medición}\\Física};
                \node[block, draw=red, fill=red!10, below=0.4cm of sim] (exp) {\textbf{Hipótesis}\\Explicación};
                
                \draw[->] (pred) -- (sim);
                \draw[->] (sim) -- (med);
                \draw[->] (med) |- (exp);
                \draw[->] (pred) |- (exp);
            \end{tikzpicture}
            
            \vspace{0.3cm}
            \begin{tcolorbox}[colback=white, colframe=gray, width=\textwidth, center]
                \centering \footnotesize
                \textbf{Cuantificación del Error (Misma Frecuencia):} \\
                \vspace{0.1cm}
                $E_G(\%) = 100 \frac{\vert G_{meas} - G_{pred}\vert}{\vert G_{pred}\vert} \quad \Delta \phi = \phi_{meas} - \phi_{pred}$
            \end{tcolorbox}
        \end{column}
    \end{columns}
\end{frame}

\begin{frame}{V8: Analogía de Redes (Eléctrico vs. Térmico)}
    \textbf{Advertencia Científica:} El modelo térmico es un \textbf{candidato} pendiente de validación en la Sesión 2.
    
    \vspace{0.2cm}
    \begin{columns}[T]
        \begin{column}{0.48\textwidth}
            \centering
            \textbf{Sistema Eléctrico RC} \\
            (Validado en Sesión 1)
            
            \vspace{0.2cm}
            {\centering $\displaystyle \tau_e \frac{dv_{out}}{dt} + v_{out} = v_{in}$ \par}
            \vspace{0.2cm}
            {\centering $\displaystyle H_e(s) = \frac{1}{1 + s\tau_e}$ \par}
            \vspace{0.2cm}
            
            \begin{itemize}
                \item Ganancia DC: $1$ (V/V) \textit{adimensional}.
                \item Dinámica: $\tau_e = RC$ [s].
            \end{itemize}
        \end{column}
        \begin{column}{0.48\textwidth}
            \centering
            \textbf{Sistema Térmico Candidato} \\
            (Planta Proyecto)
            
            \vspace{0.2cm}
            {\centering $\displaystyle C_{th} \frac{d\Delta T}{dt} + G_{th}\Delta T = P(t)$ \par}
            \vspace{0.2cm}
            {\centering $\displaystyle \frac{\Delta T(s)}{P(s)} = \frac{1/G_{th}}{1 + s \frac{C_{th}}{G_{th}}}$ \par}
            \vspace{0.2cm}
            
            \begin{itemize}
                \item Ganancia DC: $1/G_{th}$ [K/W].
                \item Dinámica: $\tau_{th} = C_{th}/G_{th}$ [s].
            \end{itemize}
        \end{column}
    \end{columns}
\end{frame}

\end{document}
